% Esse arquivo contém comandos que mudam a aparência do texto ou geram texto em algum lugar.

% === Aproxima as sections do template original ===
\titleformat*{\section}{\normalsize\bfseries}
\titleformat*{\subsection}{\normalsize\bfseries}
% === Aproxima as sections do template original ===

% === Fonte ===
\setmainfont{Times New Roman}
% === Fonte ===

% === Determina nome de figura e tabela ===
\renewcommand{\figurename}{Fig.}
\renewcommand{\tablename}{Tab.}
% === Determina nome de figura e tabela ===

% === Comandos de formatação de capa ===
\newcommand{\estudante}[1]{%
    {\large \textbf{#1}}\\
    {\setlength{\baselineskip}{0.9\baselineskip}
    \InstituicaoEstudante\\
    \EscolaEstudante\\
    \EnderecoEstudante\\
    }
}

\newcommand{\coorientadores}{%
    {\large\textbf{\CoorientadorUm%
        \ifdefempty{\CoorientadorDois}%
        {}%
        {\space\&\space\CoorientadorDois}%
    }}\\
    {\setlength{\baselineskip}{0.9\baselineskip}
        \InstituicaoCoorientadorUm\\
        \LaboratorioCoorientadorUm\\
        \EnderecoCoorientadorUm\\
    }%
    %\ifdefempty{\CoorientadorDois}%
    %{}%
    %{%
    %    {\setlength{\baselineskip}{0.9\baselineskip}
    %        \InstituicaoCoorientadorDois\\
    %        \LaboratorioCoorientadorDois\\
    %        \EnderecoCoorientadorDois\\
    %    }%
    %}%
}
\newcommand{\orientador}{%
    {\large\textbf{\Orientador}}\\
    {\setlength{\baselineskip}{0.9\baselineskip}
        \InstituicaoOrientador\\
        \LaboratorioOrientador\\
        \EnderecoOrientador\\
    }%
}
% === Comandos de formatação de capa ===